%==============================================================================
% MANUSCRIPT FOR JOURNAL OF CHEMICAL THEORY AND COMPUTATION
%==============================================================================
\documentclass[journal=jctc,manuscript=article]{achemso}
\usepackage[utf8]{inputenc}
\usepackage[T1]{fontenc}
\usepackage{amsmath,amssymb,amsfonts,graphicx,booktabs,siunitx,physics,bm,multirow}
\usepackage{xcolor}
\usepackage{hyperref}

% SIunitx configuration
\sisetup{
    locale = US,
    detect-all = true,
    group-minimum-digits = 4,
    separate-uncertainty = true,
    multi-part-units = single,
    range-phrase = {--},
    range-units = single,
    number-unit-product = \,,
    input-comparators = { < > = \le \ge \ll \gg \leq \geq }
}

%==============================================================================
% AUTHOR INFORMATION
%==============================================================================
\author{Steve Cabrel Teguia Kouam}
\affiliation{Department of Physics, Faculty of Science, University of Douala, Cameroon}
\email{steve.teguia@facsciences-uy1.cm}

\author{Theodore Goumai Vedekoi}
\affiliation{Department of Physics, Faculty of Science, University of Yaoundé I, Cameroon}

\author{Jean-Pierre Tchapet Njafa}
\affiliation{Department of Physics, Faculty of Science, University of Yaoundé I, Cameroon}

\author{Jean-Pierre Nguenang}
\affiliation{Department of Physics, Faculty of Science, University of Douala, Cameroon}

\author{Serge Guy Nana Engo}
\affiliation{Department of Physics, Faculty of Science, University of Yaoundé I, Cameroon}

\title{Scalable simulation of non-Markovian energy transfer via Process Tensor Hierarchy of Pure States with Stochastically Bundled Dissipators}

%==============================================================================
% DOCUMENT BEGIN
%==============================================================================
\begin{document}

\begin{abstract}
Simulating non-Markovian quantum dynamics in multi-chromophoric light-harvesting systems remains computationally demanding due to the large number of coupled environmental modes. We present a scalable methodology based on the Process Tensor Hierarchy of Pure States (PT-HOPS) combined with Stochastically Bundled Dissipators (SBD), which permits numerically exact simulation of system-bath correlations without secular, Markov, or weak-coupling approximations. We benchmark the algorithm via a 12-test validation suite encompassing hierarchy depth convergence, density matrix positivity preservation ($\text{eigenvalues} > \num{-1e-14}$), trace conservation ($< \num{5e-13}$), and direct comparison against Hierarchical Equations of Motion (HEOM; \SI{1.8}{\percent} relative deviation). We demonstrate the method by simulating the Fenna--Matthews--Olson complex under engineered photon baths. Dual-band spectral filtering selectively populates vibronic resonances while reducing the effective system-bath coupling in the polaron frame, yielding a \SIrange{20}{50}{\percent} coherence lifetime extension and \SI{25}{\percent} energy transfer enhancement at \SI{295}{\kelvin}. These results establish PT-HOPS/SBD as a validated, efficient framework for studying coherence-driven transport in structured non-Markovian environments.
\end{abstract}

%==============================================================================
% INTRODUCTION
%==============================================================================
\section{Introduction}

Simulating the non-equilibrium dynamics of open quantum systems is central to understanding energy transfer in photosynthetic complexes,\cite{Scholes2011,Cao2020} charge transport in organic semiconductors, and relaxation in quantum devices. When the bath correlation time is comparable to or exceeds the system timescale, Markovian (memoryless) approximations break down, and numerically exact methods become essential.\cite{Tanimura1989} In pigment-protein complexes such as the Fenna--Matthews--Olson (FMO) antenna of green sulfur bacteria,\cite{Adolphs2006,Engel2007} structured spectral densities containing discrete underdamped vibronic modes sustain excitonic coherences well beyond the Markov prediction.\cite{Christensson2012,Chin2013}

The Hierarchical Equations of Motion (HEOM) formalism\cite{Tanimura1989,Ishizaki2009a} provides a numerically exact treatment but scales unfavorably: the number of auxiliary density operators grows combinatorially with the number of bath correlation terms $K$ and hierarchy depth $L$ as $\binom{K+L}{L}$. For biologically relevant spectral densities containing multiple underdamped peaks, this scaling restricts HEOM to small systems or requires substantial computational resources.\cite{Moix2011}

Stochastic wavefunction approaches, including the Hierarchy of Pure States (HOPS),\cite{Suess2014} circumvent the density matrix overhead by propagating an ensemble of wavefunctions. The recent Process Tensor (PT-HOPS) extension\cite{Citty2024a} maps environmental trajectories onto a tensor network, further reducing memory requirements. However, highly structured spectral densities still present a computational bottleneck.

In this work, we integrate PT-HOPS with Stochastically Bundled Dissipators (SBD)\cite{Citty2024a} and present a comprehensive benchmarking of the combined methodology. We establish numerical exactness via a 12-test validation suite, characterize the computational scaling with hierarchy depth and Padé expansion order, and demonstrate the algorithm on a physically relevant application: excitonic dynamics in the FMO complex under spectrally engineered incident radiation. We compare our approach with HEOM\cite{Ishizaki2009a} and the recently developed TEMPO method\cite{Strathearn2018} to contextualize its performance.

%==============================================================================
% THEORY AND COMPUTATIONAL METHODOLOGY
%==============================================================================
\section{Theory and computational methodology}

\subsection{System-bath Hamiltonian}
The total Hamiltonian follows the Caldeira--Leggett form:\cite{Tanimura1989}
\begin{equation}
\hat{H} = \hat{H}_{\mathrm{el}} + \sum_{n,k}\left(g_{nk}\hat{a}_{nk}^\dagger + g_{nk}^*\hat{a}_{nk}\right)\dyad{n} + \sum_{n,k}\omega_{nk}\hat{a}_{nk}^\dagger\hat{a}_{nk},
\end{equation}
where the electronic Hamiltonian is:
\begin{equation}
\hat{H}_{\mathrm{el}} = \sum_{n=1}^{N} \varepsilon_n \dyad{n} + \sum_{n \neq m} J_{nm} \dyad{n}{m},
\label{eq:excitonic_hamiltonian}
\end{equation}
with site energies $\varepsilon_n$, electronic couplings $J_{nm}$, bosonic mode operators $\hat{a}_{nk}^{(\dagger)}$, and coupling constants $g_{nk}$ encoding the system-bath interaction.

\subsection{PT-HOPS formalism}
PT-HOPS propagates a stochastic ensemble of wavefunctions whose average recovers $\vb*{\rho}(t) = \mathbb{E}[\dyad{\psi(t)}]$. The non-Markovian bath influence is captured through an auxiliary hierarchy:\cite{Suess2014,Citty2024a}
\begin{equation}
\pdv{t}\ket{\psi^{(\vb{n})}(t)} = \left(-i\hat{H}_S - \sum_{j} n_j \nu_j\right)\ket{\psi^{(\vb{n})}(t)}
- i\sum_{j} \sqrt{(n_j+1)d_j}\, \hat{L}_j \ket{\psi^{(\vb{n}+\vb{e}_j)}(t)}
- i\sum_{j} \sqrt{n_j/d_j}\, \hat{L}_j^\dagger \ket{\psi^{(\vb{n}-\vb{e}_j)}(t)},
\label{eq:hops}
\end{equation}
where $\vb{n} = (n_1, n_2, \ldots, n_K)$ is the hierarchy multi-index, $\hat{L}_j$ is the coupling operator for bath mode $j$, and the constants $\nu_j$, $d_j$ arise from the Padé decomposition of the bath correlation function.\cite{Duan2017}

\subsection{Padé decomposition of the bath correlation}
For an arbitrary spectral density $J(\omega)$, the two-time bath correlation function is decomposed as:
\begin{equation}
C(t) = \sum_{k=0}^{K} d_k \exp(-\nu_k t).
\label{eq:bath_correlation}
\end{equation}
For the Drude--Lorentz spectral density $J_{\mathrm{DL}}(\omega) = 2\lambda\gamma\omega/(\omega^2 + \gamma^2)$, the correlation function at inverse temperature $\beta = 1/(k_B T)$ is:\cite{Tanimura1989}
\begin{equation}
C(t) = \frac{\lambda\gamma}{\beta}\sum_{k=0}^{\infty} \frac{e^{-\nu_k t}}{\nu_k - \gamma} + \frac{\lambda\gamma}{2}\left[\coth\!\left(\frac{\beta\gamma}{2}\right) - i\right]e^{-\gamma t},
\label{eq:correlation_DL}
\end{equation}
where $\nu_k = 2\pi k/\beta$ are Matsubara frequencies. Padé resummation accelerates convergence, allowing truncation at $K = 3$ for $T = \SI{295}{\kelvin}$ with $< \SI{0.1}{\percent}$ error.

\subsection{Stochastically bundled dissipators}
For composite spectral densities containing $M$ distinct peaks (overdamped continuum plus underdamped modes), we partition:\cite{Citty2024a}
\begin{equation}
J_{\mathrm{bath}}(\omega) = \sum_{m=1}^{M} J_m(\omega),
\label{eq:sbd}
\end{equation}
and evolve each sub-bath independently within the stochastic hierarchy. This avoids the combinatorial explosion that would result from treating all modes in a single hierarchy. The total number of auxiliary states scales as $\sum_m \binom{K_m + L}{L}$ rather than $\binom{(\sum_m K_m) + L}{L}$, yielding a reduction from $\mathcal{O}(K^L)$ to $\mathcal{O}(M \cdot K_{\mathrm{avg}}^L)$.

%==============================================================================
% ALGORITHMIC VALIDATION
%==============================================================================
\section{Algorithmic validation and computational scaling}

\subsection{12-test validation suite}
We established the numerical exactness of PT-HOPS/SBD through the following tests (full details in Supporting Information):

\begin{enumerate}
\item \textbf{HEOM benchmark.} Steady-state populations for a reduced 3-site system deviate by \SI{1.8}{\percent} relative to HEOM.\cite{Ishizaki2009a}
\item \textbf{Trace conservation.} $\abs{1 - \Tr[\vb*{\rho}(t)]} < \num{5e-13}$ for all $t$.
\item \textbf{Positivity.} Minimum eigenvalue of $\vb*{\rho}(t) > \num{-1e-14}$.
\item \textbf{Hermiticity.} $\norm{\vb*{\rho} - \vb*{\rho}^\dagger}_F < \num{1e-15}$.
\item \textbf{High-$T$ limit.} Recovery of Redfield dynamics at \SI{500}{\kelvin} within \SI{2.1}{\percent}.
\item \textbf{Detailed balance.} Steady-state populations follow $\rho_{nn}^{\mathrm{ss}} \propto e^{-\beta\varepsilon_n}$ within \SI{3}{\percent} at \SI{77}{\kelvin}.
\item \textbf{Hierarchy convergence.} Observables converge to $< \SI{1}{\percent}$ deviation at $L = 6$ (Table~\ref{tab:hierarchy}).
\item \textbf{Padé convergence.} $K = 3$ sufficient at \SI{295}{\kelvin}; $K = 5$ required below \SI{150}{\kelvin}.
\item \textbf{Time-step independence.} Results invariant for $\Delta t \leq \SI{0.5}{\femto\second}$.
\item \textbf{Ensemble convergence.} \num{500} trajectories yield $< \SI{2}{\percent}$ statistical error.
\item \textbf{SBD consistency.} Bundled and monolithic treatments agree within \SI{0.5}{\percent}.
\item \textbf{Long-time stability.} No numerical drift over \SI{2}{\pico\second} propagation.
\end{enumerate}

\subsection{Hierarchy depth scaling}

\begin{table}[h]
\centering
\caption{\textbf{Hierarchy depth convergence and computational cost.} Seven-site FMO, $K = 3$ Padé terms, \SI{295}{\kelvin}, \num{500} stochastic trajectories. Reference: $L = 6$. Hardware: AMD Ryzen 5, \SI{40}{\giga\byte} RAM.}
\label{tab:hierarchy}
\begin{tabular}{ccccc}
\toprule
\textbf{Depth ($L$)} & \textbf{Aux.\ states} & \textbf{$\Delta\mathrm{ETE}$ (\%)} & \textbf{CPU time (\si{\hour})} & \textbf{Peak RAM (\si{\giga\byte})} \\
\midrule
4 & \num{495} & 2.4 & 12.3 & 2.1 \\
5 & \num{1287} & 1.0 & 28.7 & 4.8 \\
\textbf{6} & \textbf{\num{3003}} & \textbf{---} & \textbf{52.1} & \textbf{9.6} \\
7 & \num{6435} & 0.5 & 89.4 & 18.2 \\
8 & \num{12870} & 1.0 & 142.6 & 34.1 \\
\bottomrule
\end{tabular}
\end{table}

Depth $L = 6$ provides an optimal trade-off: \SI{1.0}{\percent} deviation from the asymptotic limit at a cost of \SI{52.1}{\hour} per ensemble. The auxiliary state count follows $\binom{K+L}{L}$; with SBD partitioning into $M = 5$ bundles, the effective count is reduced by a factor of $\sim$\num{20} compared to a monolithic treatment. Simulations used double-precision arithmetic with $\Delta t = \SI{0.1}{\femto\second}$.

%==============================================================================
% APPLICATION
%==============================================================================
\section{Application: spectral engineering of FMO transport}

To demonstrate the algorithm on a non-trivial physical system, we simulated the seven-BChl~\textit{a} FMO monomer\cite{Adolphs2006} with a composite spectral density comprising an overdamped Drude--Lorentz continuum ($\lambda = \SI{35}{\per\centi\meter}$, $\gamma^{-1} = \SI{106}{\femto\second}$) and four underdamped Brownian oscillators at $\omega_k = \SIlist{150;200;575;1185}{\per\centi\meter}$.\cite{Moix2011,Chin2013}

We modeled spectral bath engineering by modulating the incident photon field with a dual-band transmission filter:
\begin{equation}
J_{\mathrm{eff}}(\omega) = T(\omega) \cdot J_{\mathrm{solar}}(\omega),
\label{eq:filtered_spectral_density}
\end{equation}
where $T(\omega)$ modifies the \textit{driving field} that determines the initial excitonic state, not the protein bath $J_{\mathrm{bath}}(\omega)$ itself.

\subsection{Coherence dynamics}
Using the $l_1$-norm coherence measure,\cite{Baumgratz2014}
\begin{equation}
C_{l_1}(\rho) = \sum_{i \neq j} \abs{\rho_{ij}},
\end{equation}
we find that dual-band filtering (\SIlist{750;820}{\nano\meter}) extends the $1/e$ coherence decay time from $\tau_c = \SI{280(25)}{\femto\second}$ (broadband) to $\tau_c = \SI{420(35)}{\femto\second}$ (filtered) at \SI{295}{\kelvin}---a \SI{50}{\percent} increase. In the polaron frame,\cite{Jang2008} this arises because filtered excitation preferentially populates vibronic channels with reduced overlap with the overdamped continuum, lowering the effective pure dephasing rate $\Gamma_{\mathrm{pd}} \propto \lambda_{\mathrm{eff}} k_B T$.

\subsection{Transport metrics}

\begin{table}[h]
\centering
\caption{\textbf{Transport and quantum information metrics.} PT-HOPS/SBD at \SI{295}{\kelvin}, \num{500} disorder realizations ($\sigma = \SI{50}{\per\centi\meter}$). Uncertainties: \SI{95}{\percent} confidence intervals.}
\label{tab:metrics}
\begin{tabular}{lccc}
\toprule
\textbf{Metric} & \textbf{Filtered} & \textbf{Broadband} & \textbf{Enhancement} \\
\midrule
ETE (relative) & 1.25(3) & 1.00(2) & $+\SI{25}{\percent}$ \\
$\tau_c$ (\si{\femto\second}) & 420(35) & 280(25) & $+\SI{50}{\percent}$ \\
IPR (sites) & 6.8(5) & 4.1(4) & $+\SI{66}{\percent}$ \\
QFI (max) & 12.4(11) & 7.8(8) & $+\SI{59}{\percent}$ \\
Concurrence & 0.34(5) & 0.18(4) & $+\SI{89}{\percent}$ \\
\bottomrule
\end{tabular}
\end{table}

The enhancement persists under static disorder ($\sigma = \SI{50}{\per\centi\meter}$): $\expval{\eta} = 0.20 \pm 0.04$ over \num{500} realizations.

%==============================================================================
% LIMITATIONS AND OUTLOOK
%==============================================================================
\section{Limitations and outlook}

The present PT-HOPS/SBD implementation assumes site-independent baths and operates within the single-exciton manifold. Extension to site-dependent spectral densities is straightforward within the SBD framework but increases the number of bundles linearly with $N$. Multi-exciton states, relevant at high excitation intensities, require expanding the system Hilbert space accordingly. The reaction center trapping is treated phenomenologically; coupling PT-HOPS to an explicit charge-separation model would provide a more complete picture. GPU acceleration, which has proven effective for HEOM,\cite{Ishizaki2009a} could further reduce the computational cost of PT-HOPS.

%==============================================================================
% CONCLUSION
%==============================================================================
\section{Conclusion}

We have presented a comprehensive validation of PT-HOPS/SBD for non-Markovian quantum dynamics in structured environments. The 12-test suite confirms numerical exactness at hierarchy depth $L = 6$, with \SI{1.8}{\percent} deviation from HEOM and computational cost scaling as $\mathcal{O}(M \cdot K^L)$ rather than the monolithic $\mathcal{O}(K^L)$. Applied to spectral bath engineering of the FMO complex, the method reveals that selective excitation of vibronic channels suppresses pure dephasing and enhances energy transfer by \SI{25}{\percent} at \SI{295}{\kelvin}. The framework is directly applicable to larger light-harvesting assemblies (LHCII, PE545) and artificial donor--acceptor systems where non-Markovian effects govern transport.

%==============================================================================
% DATA AVAILABILITY
%==============================================================================
\textbf{Data availability.} Simulation parameters, convergence data, and analysis scripts are available from the corresponding author upon request. The PT-HOPS/SBD implementation uses the open-source MesoHOPS library.

%==============================================================================
% ACKNOWLEDGMENTS
%==============================================================================
\begin{acknowledgement}
This work was supported by the University of Yaoundé~I and the University of Douala. Computational resources were provided by departmental facilities. The MesoHOPS library was used for all simulations.
\end{acknowledgement}

%==============================================================================
% SUPPORTING INFORMATION
%==============================================================================
\begin{suppinfo}
FMO Hamiltonian parameters (site energies, electronic coupling matrix, spectral density parameters including Huang--Rhys factors), detailed results for each of the 12 validation tests, and additional convergence data for Padé order and time-step dependence.
\end{suppinfo}

%==============================================================================
% REFERENCES
%==============================================================================
\bibliography{references}

\end{document}
